\documentclass{article}
\usepackage{amsmath,amssymb}
\usepackage[margin=1in]{geometry}
\usepackage[UTF8]{ctex}
\begin{document}

\section*{已知需求的随机动态定价，并附带财务里程碑相关的惩罚措施}

请选择一种非预期定价策略$\{p(t):t\in[0,T]\}$，以最大化收益。
\begin{equation}
\mathbb{E}\left[
\int_0^T p(t)\,dS(t)
-\sum_{i=1}^{m}\left[
\beta_{1,i}(\xi_i-S(\tau_i))^{+}
+\beta_{2,i}(\zeta_i-R(\tau_i))^{+}
\right]\right],
\tag{3}
\end{equation}
受…的约束
\begin{equation}
S(T)\le C\qquad\text{almost surely}.
\tag{4}
\end{equation}
对于每一个$t\in[0,T]$，非预期性意味着$p(t)$只能依赖于过去的需求观测数据和价格实际表现。这一决策就是定价策略；而销售和收入则是由该策略所引发的随机过程。论文中将$J^*$定义为满足(4)条件的所有非预期性定价策略下，式(3)的最大值。

\section*{需求与累积过程}

卖家以$C$单位的同质商品开始销售，销售时间跨度为$T$。潜在客户按照同质泊松过程到达，瞬时到达率为$\Lambda$。每位客户愿意支付$v$来购买一单位商品，该数值是从累积分布函数$F(\cdot)$中独立抽取的。在这个选定的版本中，$\Lambda$和$F(\cdot)$是决策者已知的。如果客户在时间$t$到达，并且满足条件$v\ge p(t)$，则他会购买一单位商品；否则，他不会购买任何商品，直接离开。最终的瞬时到达率就是如此。
\begin{equation}
\lambda(p(t)):=\Lambda\mathbb{P}(v\ge p(t))
=\Lambda\bar F(p(t)),
\tag{1}
\end{equation}
其中
\[
\bar F(\cdot):=1-F(\cdot).
\]
对于每一个$t\in[0,T]$，$S(t)$表示在$[0,t]$期间售出的单位数量，而$R(t)$则表示截至$t$时间的累计收入。所有时间过程都是右连续型的，并且具有左极限。用于比较里程碑的收入是指实际销售收入，不包括任何罚款；发生的罚款不会减少用于未来收入目标的收入。

\section*{里程碑数据及处罚}

在合并这些里程碑网格之前，销售目标已经明确指定了。
\[
0<t_1<t_2<\cdots<t_k\le T,
\qquad
\xi_1\le\xi_2\le\cdots\le\xi_k\le C,
\]
以及通过以下方式实现收入目标：
\[
0<t'_1<t'_2<\cdots<t'_{k'}\le T,
\qquad
\zeta_1\le\zeta_2\le\cdots\le\zeta_{k'}.
\]
这里的$\xi_i$是$t_i$计划销售的最小单位数量，而$\zeta_i$则是$t'_i$计划实现的累计收入。

该论文将这两种目标都放置在同一个网格上。
\[
\tau_0=0<\tau_1<\tau_2<\cdots<\tau_m=T,
\qquad \xi_0=0,\qquad\zeta_0=0.
\]
在这个共同的网格上，当$\tau_i$处没有销售或收入限制时，相应的目标值会被设定为零。因此，前面的单调性声明描述的是原始的目标序列，而不是对填充了零值的共同网格数组的额外单调性限制。$\tau_m=T$处的目标值是可选的；论文并不要求完全售罄或达到正的终端收入目标。

固定的处罚参数满足以下条件：
\[
\beta_{1,i}>0,\qquad\beta_{2,i}>0,
\qquad i=1,\ldots,m.
\]
在处罚定义中提到的销售和收入里程碑条件分别是：
\[
S(\tau_i)\ge\xi_i,
\qquad R(\tau_i)\ge\zeta_i,
\qquad i=1,\ldots,m.
\]
这些是可能被违反的目标，而不是在(3)--(4)中新增的硬性可行性约束条件。它们各自的不足惩罚如下：
\begin{equation}
\beta_{1,i}(\xi_i-S(\tau_i))^+
\qquad\text{and}\qquad
\beta_{2,i}(\zeta_i-R(\tau_i))^+,
\qquad i=1,\ldots,m.
\tag{2}
\end{equation}
处罚的严重程度与目标被遗漏的幅度成正比。

\section*{输入假设及其定义量}

本文提出了以下假设，这些假设在整篇论文中均成立。这些假设是对需求模型以及固定输入数据的限制，而非额外的样本路径约束条件。

\paragraph{假设 1.}
支付意愿分布在$[0,\bar p]$中有体现，其中……
\[
0<\bar p<\infty,
\]
$\bar F(\cdot)$在$[0,\bar p]$上是连续可微分的。

让$\psi(\cdot)$表示$\bar F(\cdot)$在$[0,1]$上的逆运算。通过以下方式定义收入函数：
\[
r(\lambda):=\lambda\psi(\lambda/\Lambda),
\qquad \lambda\in[0,\Lambda].
\]

\paragraph{假设 2.}
(i) $r(\lambda)$对于$\lambda\in[0,\Lambda]$来说，是严格凹的。

(ii) 通过以下方式定义能够实现收入最大化的价格及其相应的销售率：
\[
p^*=\operatorname*{arg\,max}_{p\in[0,\bar p]}p\bar F(p),
\qquad \lambda^*=\lambda(p^*).
\]
对于某些$\sigma>0$，其容量和目标都是符合要求的。
\[
\lambda^*-\sigma\ge\frac{C}{T},
\qquad
\lambda^*-\sigma\ge\frac{\xi_i}{\tau_i},
\qquad
p^*\lambda^*-\sigma\ge\frac{\zeta_i}{\tau_i},
\qquad i=1,\ldots,m.
\]

\paragraph{在假设2的解释中相关的定义。}
逆收入函数是指…
\begin{equation}
\varphi(u):=\min\{\lambda:r(\lambda)=u\},
\tag{5}
\end{equation}
以 codomain $[0,\lambda^*]$为条件。这代表了能够实现预期收入率 $u$的销售率。与产能消耗相关的价格以及各个目标值都包含在其中。
\[
p_0=\psi\!\left(\frac{C}{T\Lambda}\right),
\qquad
p_i^s=\psi\!\left(\frac{\xi_i}{\tau_i\Lambda}\right),
\qquad
p_i^r=\psi\!\left(\frac{\varphi(\zeta_i/\tau_i)}{\Lambda}\right),
\qquad i=1,\ldots,m.
\]
这里的$p_0$是指恰好在$T$时耗尽容量的价格，而不是随机定价政策的预设初始值。论文指出，假设2(ii)意味着$\lambda(p)$对于价格具有弹性。
\[
p\ge\min(p_0,p_1^s,\ldots,p_m^s,p_1^r,\ldots,p_m^r).
\]
这是需求函数的一个解释性属性，并非对决策过程$p(t)$施加的某种限制。所有数字符号都应保持原样。

\paragraph{假设 3.}
对于某些$\underline K,\bar K,\underline\kappa>0$，
\[
\underline K\le F'(p)\le\bar K
\qquad\text{on }[p^*,\bar p],
\qquad
r'\!\left(\lambda^*-\frac{\sigma}{2}\right)
\ge\underline\kappa,
\]
其中$\sigma$是假设2中使用的值。

\section*{来源范围与限制}

上述的表述和定义均来自 Besbes 和 Maglaras 的著作\emph{《动态定价与财务里程碑：反馈型政策》}，《管理科学》第58卷第9期，2012年，第2节，页码为~1717–1719页（物理PDF文件中的第3–5页）。

在(3)至(4)中打印出的政策声明中，并未单独指定$p(t)$的允许价格区间。 区间$[0,\bar p]$被明确指定用于表示支付意愿、可微性以及最大化$p^*$所定义的目标。所有数字符号都需保持原样，按原始顺序进行复制。 此处不插入额外的政策约束。 同样，式(5)中的逆函数也保持不变，包括其打印出的陪域，同时不添加任何关于空水平集的未打印约定。

这种针对特定目标的确定性表述、其可行性标准以及后来出现的足够大惩罚条件，并未适用于这种随机惩罚表述。用于后续反馈策略证明的在线辅助资料并未提供，且尚未经过验证；因此这里并未使用其中的任何公式。

\end{document}
